\FloatBarrier
\section{Calculating the Confining Potential}
\label{app:device-confinement}
\label{chap:device-models}
\label{sec:device-electrostatics}

Section~\ref{sec:quantum-dot-introduction} describes how the device
confines electrons and uses $V_{\mathrm{conf}}$ in the orbital
Hamiltonian. This appendix calculates that potential from the device
geometry, material properties, gate voltages, and background charge.
The gate and reservoir surfaces in
Fig.~\ref{fig:platform-device-section} set electrostatic boundary
conditions; solving for the electric potential and adding the material
band edge gives the electron's potential energy.


Our target is the spatial function $V_{\mathrm{conf}}(\mathbf r)$
in the electron Hamiltonian. We first specify where the electric
potential $\Phi_0(\mathbf r)$ is to be found and the surfaces on
which its boundary and interface conditions apply.
Figure~\ref{fig:device-electrostatic-regions} gives these regions
and surfaces for the reservoir--barrier--dot arrangement.
The $x,y$ coordinates lie in the material plane, and $z$ is
perpendicular to the layers.

\subsection{Define the regions and boundaries}
We first use a fixed-background electrostatic model. The gates,
electrical contact, and populated reservoir are treated as ideal
conductors. For the reservoir, this approximates strong electrostatic
screening. Each conductor has a constant potential fixed by the
external circuit.

The unknown potential is calculated in three material regions:
the silicon outside the idealized reservoir, $\mathcal D_{\mathrm{Si}}$;
the oxide, $\mathcal D_{\mathrm{ox}}$; and the surrounding air,
$\mathcal D_{\mathrm{air}}$. These regions, joined at their material
interfaces, form the solution region $\mathcal D$. The dot lies
within $\mathcal D_{\mathrm{Si}}$. Its few electrons will be included
through their explicit interaction in
Section~\ref{sec:interacting-electron-models}; here we calculate
their external potential.

The boundary of $\mathcal D$ has two parts. The conductor surfaces
form $\Gamma_{\mathrm D}$: $\mathcal S_g$ is the surface of gate $g$,
and $\mathcal S_{\mathrm R}$ is the surface of the connected
contact--reservoir region. The outer faces that limit the simulated
volume form $\Gamma_{\mathrm N}$. Thus
$\partial\mathcal D=\Gamma_{\mathrm D}\cup\Gamma_{\mathrm N}$.
Within $\mathcal D$, $\Sigma$ denotes the silicon--oxide and
oxide--air interfaces. Layer thicknesses and conductor sizes and
positions specify the three-dimensional geometry; the figure shows
an $x$--$z$ cross-section.

\begin{figure}[!t]
 \centering
 % Electrostatic solution regions and boundary types, in an x-z section.
% Coordinates are in millimetres. Conductor interiors are shaded grey.
\begingroup
\small
\setlength{\unitlength}{1mm}
\begin{picture}(150,106)
 % Unknown-potential regions.
 \put(8,28){\color{blue!7}\rule{135mm}{31mm}}
 \put(8,59){\color{black!7}\rule{135mm}{12mm}}
 \put(8,71){\color{blue!2}\rule{135mm}{25mm}}

 % Ideal conductors: the contact and reservoir form one connected region.
 \put(17,49){\color{black!20}\rule{16mm}{34mm}}
 \put(33,49){\color{black!20}\rule{42mm}{7mm}}
 \put(41,71){\color{black!15}\rule{37mm}{12mm}}
 \put(87,71){\color{black!15}\rule{19mm}{12mm}}
 \put(115,71){\color{black!15}\rule{23mm}{12mm}}

 % Sigma: dotted material interfaces stop at conductor surfaces.
 \color{green!40!black}
 \multiput(8,59)(1,0){10}{\circle*{0.35}}
 \multiput(33,59)(1,0){111}{\circle*{0.35}}
 \multiput(8,71)(1,0){10}{\circle*{0.35}}
 \multiput(33,71)(1,0){9}{\circle*{0.35}}
 \multiput(78,71)(1,0){10}{\circle*{0.35}}
 \multiput(106,71)(1,0){10}{\circle*{0.35}}
 \multiput(138,71)(1,0){6}{\circle*{0.35}}

 % Gamma_D: solid conductor surfaces facing the solution region.
 \color{red!65!black}
 \linethickness{0.8pt}
 \put(17,49){\line(0,1){34}}
 \put(17,83){\line(1,0){16}}
 \put(33,56){\line(0,1){27}}
 \put(33,56){\line(1,0){42}}
 \put(75,49){\line(0,1){7}}
 \put(17,49){\line(1,0){58}}
 \put(41,71){\line(1,0){37}}
 \put(41,83){\line(1,0){37}}
 \put(41,71){\line(0,1){12}}
 \put(78,71){\line(0,1){12}}
 \put(87,71){\line(1,0){19}}
 \put(87,83){\line(1,0){19}}
 \put(87,71){\line(0,1){12}}
 \put(106,71){\line(0,1){12}}
 \put(115,71){\line(1,0){23}}
 \put(115,83){\line(1,0){23}}
 \put(115,71){\line(0,1){12}}
 \put(138,71){\line(0,1){12}}

 % Gamma_N: the outer boundary of the finite computational volume.
 \color{blue!70!black}
 \linethickness{0.55pt}
 \put(8,28){\dashbox{1.5}(135,68){}}
 \put(8,99){\makebox(90,5)[l]{$\Gamma_{\mathrm N}$: outer boundary}}

 % All text uses the same font size.
 \color{black}
 \put(40,87){\makebox(43,5){$\mathcal D_{\mathrm{air}}$: air}}
 \put(78,62){\makebox(60,5){$\mathcal D_{\mathrm{ox}}$: oxide}}
 \put(39,34){\makebox(74,5){$\mathcal D_{\mathrm{Si}}$: silicon}}
 \put(17,64){\makebox(16,7){contact}}
 \put(33,49){\makebox(42,7){reservoir ($\Phi_{\mathrm R}$)}}
 \put(41,71){\makebox(37,12){\shortstack{accumulation\\gate}}}
 \put(87,71){\makebox(19,12){\shortstack{barrier\\gate}}}
 \put(115,71){\makebox(23,12){\shortstack{plunger\\gate}}}
 \put(116,49){\color{orange!70!black}\makebox(22,6){dot}}

 % Surface labels point to the boundary, not the conductor interior.
 \color{red!65!black}
 \put(87,87){\makebox(19,5){$\mathcal S_g$}}
 \put(96.5,87){\vector(0,-1){4}}
 \put(46,40){\makebox(18,5){$\mathcal S_{\mathrm R}$}}
 \put(55,45){\vector(0,1){4}}
 \color{green!40!black}
 \put(78.5,75){\makebox(8,5){$\Sigma$}}
 \put(82.5,75){\vector(0,-1){4}}
 \put(82,50){\makebox(8,5){$\Sigma$}}
 \put(86,55){\vector(0,1){4}}

 % Coordinate axes within the silicon region.
 \color{black}
 \linethickness{0.4pt}
 \put(14,33){\vector(1,0){13}}
 \put(14,33){\vector(0,1){11}}
 \put(28,30.5){\makebox(5,5){$x$}}
 \put(11.5,44){\makebox(5,5){$z$}}

 % Legend: line style, colour, and symbols all identify each boundary type.
 \color{red!65!black}
 \linethickness{0.8pt}
 \put(12,18){\line(1,0){14}}
 \put(30,15){\makebox(110,6)[l]{$\Gamma_{\mathrm D}$: conductor surfaces}}
 \color{blue!70!black}
 \linethickness{0.55pt}
 \multiput(12,10)(3,0){5}{\line(1,0){2}}
 \put(30,7){\makebox(110,6)[l]{$\Gamma_{\mathrm N}$: outer boundary}}
 \color{green!40!black}
 \multiput(12,2)(1,0){15}{\circle*{0.35}}
 \put(30,-1){\makebox(110,6)[l]{$\Sigma$: material interfaces}}
\end{picture}
\endgroup

 \caption{Regions and boundaries of the electrostatic model.
 The labeled silicon, oxide, and air regions form $\mathcal D$.
 Solid red lines mark conductor surfaces $\Gamma_{\mathrm D}$;
 dashed blue lines mark the outer boundary $\Gamma_{\mathrm N}$;
 dotted green lines mark material interfaces $\Sigma$.
 The $y$ direction is perpendicular to the drawing; the undrawn
 front and back outer faces also belong to $\Gamma_{\mathrm N}$.
 The schematic is not to scale.}
 \label{fig:device-electrostatic-regions}
\end{figure}

\subsection{Specify material data and boundary conditions}
Electric fields displace bound charges in a material. Each
material's permittivity $\epsilon_{\mathrm d}$ describes the
resulting polarization. Its value in each part of the device
defines the function $\epsilon_{\mathrm d}(\mathbf r)$ over
$\mathcal D$. The material conduction-band edge
$E_{\mathrm c}^{0}(\mathbf r)$ is another input, incorporating
band offsets and strain with energies measured from a common reference.
Band offsets are energy steps at material boundaries. Strain is an
elastic deformation of the crystal that changes its band energies;
its effect is included in the supplied material profile.

The remaining source is a prescribed background charge density
$\rho_{\mathrm{bg}}(\mathbf r)$, measured in charge per unit
volume. A charge-free baseline sets $\rho_{\mathrm{bg}}=0$
throughout $\mathcal D$; a specified impurity-charge profile can
be added here. The reservoir's screening enters through its
conducting boundary. Its surface charge will adjust to maintain
the prescribed potential.

On the gate surfaces $\mathcal S_g\subset\Gamma_{\mathrm D}$, impose
\begin{equation}
 \left.\Phi_0\right|_{\mathcal S_g}
       =V_g+\Phi_g^{\mathrm{off}},
 \label{eq:device-gate-boundary}
\end{equation}
where $V_g$ is the applied voltage and
$\Phi_g^{\mathrm{off}}$ is a specified built-in offset from the
material and voltage-reference conventions.
On $\mathcal S_{\mathrm R}$, the contact and reservoir have the
prescribed potential $\Phi_{\mathrm R}$. This is an electrostatic boundary
value; the reservoir electrochemical potential $\mu_{\mathrm R}$
sets the occupation of electron states.

These prescribed conductor values define a function $\Phi_{\mathrm D}$
on $\Gamma_{\mathrm D}$. All potentials refer to the same reference
conductor.

On the outer faces $\Gamma_{\mathrm N}$, take zero normal electric
displacement:
$\mathbf n\cdot\epsilon_{\mathrm d}\nabla\Phi_0=0$,
where $\mathbf n$ is the outward normal. The outer faces are placed
far enough away that moving them farther changes the potential
near the dot negligibly.

Across the material interfaces $\Sigma$, the solutions in neighboring
materials must join. We assume these interfaces carry no additional charge,
such as trapped charge; material polarization is already included
in $\epsilon_{\mathrm d}$. The potential and normal electric
displacement are then continuous across each interface.

\subsection{Solve the resulting boundary-value problem}
The electric field is $\mathbf E=-\nabla\Phi_0$. Substitution
into Gauss's law gives the problem for the unknown $\Phi_0$:
\begin{equation}
 \begin{aligned}
 -\nabla\cdot[\epsilon_{\mathrm d}\nabla\Phi_0]
       &=\rho_{\mathrm{bg}}
            &&\text{in }\mathcal D\setminus\Sigma,\\
 \Phi_0&=\Phi_{\mathrm D}
            &&\text{on }\Gamma_{\mathrm D},\\
 \mathbf n\cdot\epsilon_{\mathrm d}\nabla\Phi_0&=0
            &&\text{on }\Gamma_{\mathrm N},\\
 \Phi_{0,1}&=\Phi_{0,2}
            &&\text{on }\Sigma,\\
 \epsilon_{\mathrm d,1}\,\partial_n\Phi_{0,1}
       &=\epsilon_{\mathrm d,2}\,\partial_n\Phi_{0,2}
            &&\text{on }\Sigma.
 \end{aligned}
\label{eq:device-poisson}
\end{equation}
The subscripts $1,2$ identify the two sides of each material
interface; $\partial_n$ uses the same normal direction on both sides.
The geometry fixes the domain and interfaces, material data
fix the coefficients, and the preceding steps fix the source
and every boundary value. With positive permittivity and a
nonempty fixed-potential boundary on each connected component,
this is a closed linear problem~\cite{QTCADPoisson}.

\subsection{Convert the electric potential to electron energy}
The resulting electron potential energy is
\begin{equation}
 V_{\mathrm{conf}}(\mathbf r)
       =E_{\mathrm c}^{0}(\mathbf r)-e\Phi_0(\mathbf r).
 \label{eq:device-confinement-energy}
\end{equation}
Appendix~\ref{sec:device-bound-orbitals} inserts this function into the
effective-mass Schr\"odinger equation to obtain dot orbitals and energies.
For several retained valleys, indexed by $\nu$, use their
respective material profiles:
$V_{\mathrm{conf},\nu}=E_{\mathrm c,\nu}^{0}-e\Phi_0$.

This construction uses a fixed background charge and ideal
reservoir screening. To resolve the reservoir's finite charge
response, include its interior in $\mathcal D$ and supply a
carrier-density model with its material parameters, temperature,
and electrochemical potential. The charge then depends on
$\Phi_0$, and the source and potential are solved
self-consistently~\cite{QTCADPoisson}. The corresponding treatment
of quantum-confined electrons is a many-electron mean-field
approximation, developed in
Section~\ref{sec:device-schrodinger-poisson}.
